\documentclass{article}%
\usepackage[T1]{fontenc}%
\usepackage[utf8]{inputenc}%
\usepackage{lmodern}%
\usepackage{textcomp}%
\usepackage{lastpage}%
\usepackage{geometry}%
\geometry{margin=2.2cm}%
%
\usepackage{amsmath}%
\usepackage{amssymb}%
\usepackage{booktabs}%
\usepackage{longtable}%
\usepackage{array}%
\usepackage{hyperref}%
\hypersetup{hidelinks}%
\title{IME 2027 {-} Matemática (claude\_haiku\_5\_5) x Poliedro\\[4pt]\large Comparação com o gabarito Poliedro\thanks{Gerado com gpt-resolve: \url{https://github.com/lgabs/gpt-resolve}}}%
\date{09/11/2026}%
%
\begin{document}%
\normalsize%
\maketitle%
\section*{Resumo}\textbf{Nota média: 10,0/10} (12 de 12 itens com a resolução correta; cada item vale o mesmo dentro da questão).\par\medskip%
\begin{longtable}{@{}c c p{0.78\textwidth}@{}}\toprule Questão & Nota & Resumo \\ \midrule
1 & 10,0 & Questão 1: a resolução e o gabarito Poliedro estão de acordo, e os dois estão certos: a área de P ? I é 12$\surd$3 - 18 $\approx$ 2,785. Reli o enunciado no PDF original e ele confere com a transcrição. O cálculo foi refeito de forma exata e numérica (recorte dos polígonos em python3). A resolução é mais detalhada que o gabarito: prova que a interseção é um dodecágono regular e confere o resultado de três jeitos (pelas pontas, numericamente e por cotas de grandeza). Não há erro que afete a nota. \\
2 & 10,0 & A resolução e o gabarito Poliedro estão certos e dão a mesma resposta, P = 2/5. A transcrição do enunciado confere com o PDF original. A resolução descarta z = 1 corretamente (espaço amostral de 5 raízes), testa cada raiz e confirma o resultado pela identidade |z-3/2|$^2$ = 13/4 - 3Re(z). Também explica que z$_{1}$ e z$_{5}$ ficam exatamente sobre o bordo do disco e que o bordo está incluído. Não há divergência nem erro. \\
3 & 10,0 & A resolução e o gabarito Poliedro dão a mesma resposta, x = 7, e ela está certa. Para x inteiro positivo, S(x) = 1 - 1/(2$\cdot$3\^{}(x-1)) $\geq$ 0,999 vale se e só se 3\^{}(x-1) $\geq$ 500, ou seja, x $\geq$ 7. Os inteiros negativos dão S = 1/2. A resolução é mais completa que o gabarito: descarta os inteiros negativos, exclui x = 0 e ainda encontra o conjunto solução real (-1/6, 0) ? (6, +$\infty$). O gabarito supõe x > 0 sem justificar. \\
4 & 10,0 & A resolução está correta e coincide com o gabarito Poliedro: o produto vale 1345. O método é o mesmo nos dois. Tirar os y\_i com a identidade ?(t+y\_i)=t$^3$-4t+2 e reduzir x$^3$ módulo x$^2$-7x-3 leva a (48x$_{1}$+23)(48x$_{2}$+23)=1345. Conferi a conta de forma independente, numericamente com raízes calculadas, e deu 1345,0. Não há erros nem pontos duvidosos. \\
5 & 10,0 & A resolução e o gabarito do Poliedro concordam. O dado impresso 5540 (confirmado no PDF oficial) é quase certamente um erro de digitação de 5440, e com 5440 a soma pedida é 10912, com x\_k = 2\^{}(k+4). A resolução também resolve o enunciado como foi impresso: a soma é $\approx$ 11196,5, um valor irracional, com p $\approx$ 2,5116 a raiz real única de 128p$^3$ - 256p - 1385 = 0. Conferi todas as contas numericamente: unicidade da raiz, ausência de raiz racional, forma fechada com cosh, fatoração para 5440 e as somas. Não há erro matemático. O único reparo ao gabarito é que a frase 'não tem solução em forma fechada' é imprecisa: a solução existe, só não é um número simples. A nota final depende da decisão da banca sobre o erro do enunciado. \\
6 & 10,0 & A resolução está certa e bate com o gabarito Poliedro: P = 1/3. Conferi o enunciado no PDF original (página 2) e ele corresponde à transcrição. O critério é o mesmo nos dois: o comprimento do ciclo de 1 divide 9, então L $\in$ \{1,3,9\}. Cada comprimento de ciclo aparece em 8! permutações. Refiz a conta do gabarito para L=3 (C(8,2)$\cdot$2$\cdot$6! = 8!) e enumerei as 9! permutações: 120960/362880 = 1/3. A argumentação da resolução está completa. Não há divergência. \\
7 & 10,0 & A resolução está correta e completa. Pelos gabaritos dos dois lados, tr(A$\cdot$B) = C(m+n+2,n+1) = C(2026,1001), resultado igual ao do gabarito Poliedro e conferido numericamente (o traço é $\equiv$ 0 mod 1013, mas não mod 1013$^2$). O argumento de divisibilidade da resolução está certo: cancela 1025! e usa o fator 2026 = 2$\cdot$1013 em 1026$\cdot$\ldots{}$\cdot$2026, junto com 1013 ? 1001!. O gabarito chega à mesma conclusão, mas sua justificativa final tem um passo falso: afirma que 1013 não divide 1025!$\cdot$1001!, quando 1013 | 1025!. \\
8 & 10,0 & A resolução e o gabarito Poliedro coincidem. Os dois demonstram S\_n = sen(n$\alpha$)/(sen(2$\alpha$)cos((n+1)$\alpha$)) pelo mesmo caminho: prostaférese, telescopagem com tangentes e diferença de tangentes. A resolução está correta e completa: justifica cada divisão pelas hipóteses, trata o caso n = 0 e mostra que as hipóteses são necessárias. As verificações numéricas que ela traz foram reconferidas em Python e estão certas. Não encontrei erro em nenhum dos dois. \\
9 & 10,0 & A resolução está correta e bate com o gabarito Poliedro: o lugar geométrico de M é 4x$^2$/49 + y$^2$ = 1 e a interseção com x = $\pm$2$\surd$3 dá os pontos ($\pm$2$\surd$3, $\pm$1/7). Conferi a transcrição contra o PDF original (pág. 2) e refiz as contas numericamente. A resolução é até mais cuidadosa que o gabarito, porque exclui os casos degenerados P=($\pm$4,0). Não há erro em nenhum dos dois. \\
10 & 10,0 & A resolução e o gabarito concordam nos dois itens: (a) V\_max = 192$\surd$3 cm$^3$ $\approx$ 332,55 cm$^3$; (b) o prisma é um cubo de aresta 4$\surd$3 cm, com altura e lado da base iguais a 4$\surd$3 cm. Refiz as contas e conferi numericamente; ambos estão corretos. A resolução é mais rigorosa que o gabarito: prova a cota para base quadrangular qualquer e bases em qualquer posição, e mostra que o máximo é atingido e é único. O gabarito supõe, sem justificar, base retangular e os 8 vértices na superfície, e traz um erro de digitação (4$\surd$(3 cm)) que não altera nada. \\
\bottomrule\end{longtable}%
\section*{Principais pontos de dúvida}%
\begin{itemize}%
\item%
\textbf{Q5} \textit{(statement)} O caderno oficial (página 2 do PDF) imprime x2+x4+x6+x8 = 5540. Com esse dado, nenhuma razão inteira serve e a resposta é um número irracional ($\approx$ 11196,5, raiz de uma cúbica sem raiz racional). O valor 5440 (d = 1, resposta 10912) é quase certamente o pretendido. A resposta que vale nota depende da decisão da banca: aceitar 10912, aceitar $\approx$ 11196,5, aceitar as duas ou anular a questão, como sugere o Poliedro. A resolução apresenta as duas respostas, o que a protege em qualquer cenário.%
\item%
\textbf{Q7} \textit{(key)} A justificativa final do gabarito Poliedro tem um passo falso: diz que 'o primo 1013 não divide 1025!$\cdot$1001!', mas 1013 $\leq$ 1025, então 1013 divide 1025!. Na verdade 2026! tem 1013$^2$, 1025! tem 1013$^{1}$ e 1001! não tem fator 1013, de modo que v\_1013(C(2026,1001)) = 1. A conclusão continua certa, mas uma prova escrita como no gabarito seria logicamente inválida nesse ponto. Vale avisar o professor para não usar esse argumento como padrão de correção.%
\end{itemize}%
\section*{Questão 1}%
\href{https://poliedroresolve.sistemapoliedro.com.br/vestibulares/ime/2027/ime-2-fase-dia-1-26-10-2026-2-fase-dia-1/dissertativa-1-matematica-2-fase-dia-1-ime-2027}{Resolução do gabarito}\par\smallskip%
\begin{longtable}{@{}p{0.06\textwidth} p{0.33\textwidth} p{0.33\textwidth} p{0.1\textwidth} p{0.08\textwidth}@{}}\toprule Item & Resolução & Gabarito & Veredito & Certo \\ \midrule
Área de P ? I & 12$\surd$3 - 18 $\approx$ 2,78 u.a. & 12$\surd$3 - 18 ($\approx$ 2,785) & igual & ambos \\
\bottomrule\end{longtable}%
\begin{itemize}%
\item%
\textbf{Área de P ? I)} Os dois chegam ao mesmo valor exato pelo mesmo caminho: cada hexágono tem área 3$\surd$3/2; a interseção é o dodecágono regular de apótema $\surd$3/2, com área 12$\cdot$(3/4)$\cdot$tg15$^\circ$ = 18 - 9$\surd$3; por inclusão{-}exclusão, a união vale 3$\surd$3 - (18 - 9$\surd$3) = 12$\surd$3 - 18. Conferi com python3, recortando os hexágonos por Sutherland--Hodgman: interseção 2,4115427 e união 2,7846097, iguais aos valores exatos. Também conferi a verificação da resolução feita pelas 6 pontas: cada ponta tem área (2$\surd$3-3)(1-$\surd$3/2)/2 = (7$\surd$3-12)/4 $\approx$ 0,0311, e 3$\surd$3/2 + 6$\cdot$(7$\surd$3-12)/4 = 12$\surd$3 - 18. A diferença entre $\approx$2,78 e $\approx$2,785 é só de arredondamento.%
\end{itemize}%
\paragraph{Pontos de dúvida}%
\begin{itemize}%
\item%
\textit{(statement)} O enunciado não diz se "hexágono" é a região poligonal (cheia) ou só o contorno. A resolução e o gabarito tratam como região cheia, que é o natural, já que se pede "área da região coberta". O resultado não muda se os índices forem de 1 a 12, como a própria resolução observa.%
\item%
\textit{(solution)} No passo 1, a frase "Se fossem iguais, a união teria área 3$\surd$3/2, o que não ocorre" sai um pouco do argumento e não justifica nada. Tirando isso, a demonstração de que P?I é o dodecágono regular, por semiplanos com normais a cada 30$^\circ$ e distância $\surd$3/2, está correta e completa.%
\end{itemize}%
\section*{Questão 2}%
\href{https://poliedroresolve.sistemapoliedro.com.br/vestibulares/ime/2027/ime-2-fase-dia-1-26-10-2026-2-fase-dia-1/dissertativa-2-matematica-2-fase-dia-1-ime-2027}{Resolução do gabarito}\par\smallskip%
\begin{longtable}{@{}p{0.06\textwidth} p{0.33\textwidth} p{0.33\textwidth} p{0.1\textwidth} p{0.08\textwidth}@{}}\toprule Item & Resolução & Gabarito & Veredito & Certo \\ \midrule
Q2 (probabilidade) & P = 2/5 = 0,4 (só z = 1/2 $\pm$ i$\surd$3/2 satisfazem, sobre o bordo do disco) & P = 2/5 & igual & ambos \\
\bottomrule\end{longtable}%
\begin{itemize}%
\item%
\textbf{Q2 (probabilidade))} Refiz a conta. Como (z-1)P(z) = z$^{6}$-1 e P(1) = 6 $\neq$ 0, as raízes são cis(k$\pi$/3), k = 1..5. Para |z| = 1, |z-3/2|$^2$ = 13/4 - 3cos $\theta$, o que dá 7/4 para k = 1,5 (igualdade, bordo incluído porque a desigualdade é $\leq$), 19/4 para k = 2,4 e 25/4 para k = 3. São 2 casos favoráveis em 5, logo P = 2/5. Conferi numericamente em python3: |P(z\_k)| \textasciitilde{} 1e{-}15 e |z-3/2|$^2$ = 1,75; 4,75; 6,25; 4,75; 1,75. Os valores coincidem com os do gabarito ($\surd$7/2, $\surd$19/2, 5/2).%
\end{itemize}%
\paragraph{Pontos de dúvida}%
\begin{itemize}%
\item%
\textit{(statement)} O enunciado não diz que as 5 raízes são equiprováveis. É a leitura natural de 'escolher ao acaso uma raiz qualquer', e tanto a resolução quanto o gabarito a adotam.%
\item%
\textit{(statement)} O resultado depende de o bordo estar incluído ($\leq$). Conferi no PDF original: o sinal é mesmo $\leq$. Com desigualdade estrita a probabilidade seria 0, como a própria resolução observa.%
\end{itemize}%
\section*{Questão 3}%
\href{https://poliedroresolve.sistemapoliedro.com.br/vestibulares/ime/2027/ime-2-fase-dia-1-26-10-2026-2-fase-dia-1/dissertativa-3-matematica-2-fase-dia-1-ime-2027}{Resolução do gabarito}\par\smallskip%
\begin{longtable}{@{}p{0.06\textwidth} p{0.33\textwidth} p{0.33\textwidth} p{0.1\textwidth} p{0.08\textwidth}@{}}\toprule Item & Resolução & Gabarito & Veredito & Certo \\ \midrule
Questão 3 (menor inteiro x com f(x)+f({-}1/x) $\geq$ 0,999) & x = 7 (S(7) = 1457/1458 $\approx$ 0,99931; S(6) = 485/486 $\approx$ 0,99794) & x = 7 & igual & ambos \\
\bottomrule\end{longtable}%
\begin{itemize}%
\item%
\textbf{Questão 3 (menor inteiro x com f(x)+f({-}1/x) $\geq$ 0,999))} Refiz a conta. Para x $\geq$ 0 todos os g\_n valem 1, então f(x) = 1/2. Para x inteiro positivo k, g\_n({-}1/k) = 1 exatamente quando n < k, o que dá S(k) = 1 - 1/(2$\cdot$3\^{}(k-1)). A condição S(k) $\geq$ 0,999 equivale a 3\^{}(k-1) $\geq$ 500, logo k $\geq$ 7. Para inteiros x $\leq$ -1, f(x) = 0 e S = 1/2. Conferi em python com frações exatas: S(6) $\approx$ 0,99794 e S(7) $\approx$ 0,99931. O enunciado do PDF bate com a transcrição.%
\end{itemize}%
\paragraph{Pontos de dúvida}%
\begin{itemize}%
\item%
\textit{(key)} O gabarito parte de x $\in$ $\mathbb{Z}$+ sem justificar. Falta descartar os inteiros negativos (para x $\leq$ -1, f(x) = 0 e f(-1/x) = 1/2, então S = 1/2 < 0,999) e excluir x = 0, que torna f(-1/x) indefinida. Como o enunciado pede o menor inteiro, numa correção rigorosa essa omissão poderia custar parte da nota. A resolução trata esses dois casos explicitamente.%
\item%
\textit{(statement)} O enunciado não diz que x = 0 está fora do domínio da expressão, mas isso decorre de -1/x. Também não fica claro se "menor valor inteiro" admite negativos. Isso não muda a resposta, porque nenhum inteiro $\leq$ 0 satisfaz a desigualdade.%
\end{itemize}%
\section*{Questão 4}%
\href{https://poliedroresolve.sistemapoliedro.com.br/vestibulares/ime/2027/ime-2-fase-dia-1-26-10-2026-2-fase-dia-1/dissertativa-4-matematica-2-fase-dia-1-ime-2027}{Resolução do gabarito}\par\smallskip%
\begin{longtable}{@{}p{0.06\textwidth} p{0.33\textwidth} p{0.33\textwidth} p{0.1\textwidth} p{0.08\textwidth}@{}}\toprule Item & Resolução & Gabarito & Veredito & Certo \\ \midrule
Questão 4 (valor do produto) & 1345 & 1345 & igual & ambos \\
\bottomrule\end{longtable}%
\begin{itemize}%
\item%
\textbf{Questão 4 (valor do produto))} Conferi o enunciado no PDF original (página 1) e ele coincide com a transcrição. A resolução e o gabarito seguem o mesmo caminho. Primeiro usam ?(y\_i+t) = -f(-t) = t$^3$-4t+2. Depois reduzem x$^3$ = 52x+21 módulo x$^2$-7x-3, de modo que cada fator vira 48x\_j+23. Por fim, com Girard (x$_{1}$+x$_{2}$=7, x$_{1}$x$_{2}$=-3), chegam a 2304$\cdot$(-3)+1104$\cdot$7+529 = 1345. No gabarito os fatores aparecem com o sinal trocado, (-48x\_j-23), mas como são dois o resultado é o mesmo. Conferi em Python com as raízes reais de y$^3$-4y-2 (2,2143; -0,5392; -1,6751) e x = (7$\pm$$\surd$61)/2. O produto numérico deu 1344,99999\ldots{}, e cada produto parcial bate com 48x\_j+23 (378,446 e 3,554). Também estão certos o discriminante (148), a cota |y\_i|<3 e o argumento de sinal.%
\end{itemize}%
\section*{Questão 5}%
\href{https://poliedroresolve.sistemapoliedro.com.br/vestibulares/ime/2027/ime-2-fase-dia-1-26-10-2026-2-fase-dia-1/dissertativa-5-matematica-2-fase-dia-1-ime-2027}{Resolução do gabarito}\par\smallskip%
\begin{longtable}{@{}p{0.06\textwidth} p{0.33\textwidth} p{0.33\textwidth} p{0.1\textwidth} p{0.08\textwidth}@{}}\toprule Item & Resolução & Gabarito & Veredito & Certo \\ \midrule
Questão 5: x1+x3+x5+x7+x9 & Com o valor impresso (x2+x4+x6+x8 = 5540), a soma é O = 5540p - 512p$^2$ + 512 $\approx$ 11196,5. Aqui p $\approx$ 2,5116 é a única raiz real de 128p$^3$ - 256p - 1385 = 0, e O é irracional. Com o valor provavelmente pretendido, 5440, a soma é 10912, com x\_k = 2\^{}(k+4) e razão q = 2. & Com 5540 a questão não tem solução em forma fechada. O gabarito supõe erro de digitação (5440), o que dá x1+x3+x5+x7+x9 = 10912 com q = 2, e sugere anular a questão. & igual & ambos \\
\bottomrule\end{longtable}%
\begin{itemize}%
\item%
\textbf{Questão 5: x1+x3+x5+x7+x9)} Conferi na página 2 do PDF oficial: o dado impresso é mesmo 5540. A resolução e o gabarito concordam no essencial: x5 = 512, a equação u$^3$ - 2u = 1385/128 com u = q + 1/q, o diagnóstico de erro de digitação e a resposta 10912 para 5440. Refiz as contas em Python:\newline%
? raiz real única p = 2,5116001\ldots{}; t = 2,0154274; d = 1,0110858;\newline%
? E = 5540,000 e O = 11196,4994, que também sai da identidade O = pE - 512(p$^2$ - 1);\newline%
? a forma fechada (2$\surd$6/3)$\cdot$cosh((1/3)$\cdot$arccosh(4155$\surd$6/1024)) dá o mesmo p;\newline%
? f(5/2) = -25 e f(277/128) = -10514643/16384, então não há raiz racional e O é irracional;\newline%
? com d inteiro, E vale 2048, 5440, 34952, \ldots{};\newline%
? com 5440: 8p$^3$ - 16p - 85 = (2p - 5)(4p$^2$ + 10p + 17), logo p = 5/2, t = 2 e O = 32 + 128 + 512 + 2048 + 8192 = 10912.\newline%
A resolução vai além do gabarito: também resolve o enunciado como foi impresso. Esse valor está correto ($\approx$ 11196,5) e não contradiz o gabarito, que só diz que não há forma fechada simples.%
\end{itemize}%
\paragraph{Pontos de dúvida}%
\begin{itemize}%
\item%
\textit{(statement)} O caderno oficial (página 2 do PDF) imprime x2+x4+x6+x8 = 5540. Com esse dado, nenhuma razão inteira serve e a resposta é um número irracional ($\approx$ 11196,5, raiz de uma cúbica sem raiz racional). O valor 5440 (d = 1, resposta 10912) é quase certamente o pretendido. A resposta que vale nota depende da decisão da banca: aceitar 10912, aceitar $\approx$ 11196,5, aceitar as duas ou anular a questão, como sugere o Poliedro. A resolução apresenta as duas respostas, o que a protege em qualquer cenário.%
\item%
\textit{(key)} O gabarito diz que, com 5540, a questão 'não tem solução em forma fechada'. Isso é impreciso. A solução existe, é única e tem forma fechada: Cardano dá uma raiz real em radicais reais (discriminante (1385/256)$^2$ - 8/27 > 0), e a resolução escreve a forma com cosh/arccosh. O que falta é um valor 'bonito' ou inteiro. A própria resolução resumida do gabarito fala em 'sem forma fechada simples', que é a formulação correta. Isso pesa num eventual argumento de anulação: a questão não é impossível, apenas tem dado inconsistente com a intenção.%
\item%
\textit{(solution)} No Passo 10, a resolução diz que, 'em ambos os casos', x9 vale cerca de 8,4$\cdot$10$^3$. Isso só vale para o caso 5540 (x9 $\approx$ 8447,7). No caso 5440, x9 = 8192 $\approx$ 8,2$\cdot$10$^3$. É uma imprecisão de ordem de grandeza sem efeito na resposta.%
\end{itemize}%
\section*{Questão 6}%
\href{https://poliedroresolve.sistemapoliedro.com.br/vestibulares/ime/2027/ime-2-fase-dia-1-26-10-2026-2-fase-dia-1/dissertativa-6-matematica-2-fase-dia-1-ime-2027}{Resolução do gabarito}\par\smallskip%
\begin{longtable}{@{}p{0.06\textwidth} p{0.33\textwidth} p{0.33\textwidth} p{0.1\textwidth} p{0.08\textwidth}@{}}\toprule Item & Resolução & Gabarito & Veredito & Certo \\ \midrule
Probabilidade de f\^{}(9)(1) = 1 & P = 1/3 & P = 1/3 & igual & ambos \\
\bottomrule\end{longtable}%
\begin{itemize}%
\item%
\textbf{Probabilidade de f\^{}(9)(1) = 1)} Os dois resolvem do mesmo jeito: f\^{}(9)(1)=1 vale exatamente quando o comprimento L do ciclo que contém 1 divide 9, ou seja, L $\in$ \{1,3,9\}. A resolução mostra que N\_L = {[}8!/(9-L)!{]}$\cdot$(9-L)! = 8! para todo L. O gabarito conta caso a caso: para L=3, C(8,2)$\cdot$2$\cdot$6! = 28$\cdot$2$\cdot$720 = 40320 = 8!, o que está certo. Então P = 3$\cdot$8!/9! = 1/3. Conferi com uma enumeração completa em Python: 120960 casos favoráveis em 362880, razão exatamente 1/3. A resolução também está completa: o Passo 1 usa a casa dos pombos e a injetividade para obter o ciclo, e o Passo 2 prova a equivalência f\^{}(m)(1)=1 $\Leftrightarrow$ L | m.%
\end{itemize}%
\paragraph{Pontos de dúvida}%
\begin{itemize}%
\item%
\textit{(solution)} A 'Conferência' cita uma enumeração completa das 9! bijeções para confirmar N = 120960. Isso não é possível numa prova feita à mão. É só uma checagem extra e a argumentação não depende dela, mas numa versão para o aluno talvez convenha tirar ou deixar claro que foi feita por computador.%
\end{itemize}%
\section*{Questão 7}%
\href{https://poliedroresolve.sistemapoliedro.com.br/vestibulares/ime/2027/ime-2-fase-dia-1-26-10-2026-2-fase-dia-1/dissertativa-7-matematica-2-fase-dia-1-ime-2027}{Resolução do gabarito}\par\smallskip%
\begin{longtable}{@{}p{0.06\textwidth} p{0.33\textwidth} p{0.33\textwidth} p{0.1\textwidth} p{0.08\textwidth}@{}}\toprule Item & Resolução & Gabarito & Veredito & Certo \\ \midrule
Demonstração (traço de A$\cdot$B divisível por 1013 para m = 2\^{}10, n = 10\^{}3) & tr(A$\cdot$B) = C(m+n+2, n+1), obtido somando a1j + a2j por Stifel e aplicando Vandermonde duas vezes. Com m = 1024 e n = 1000, tr = C(2026,1001). Como 1001!$\cdot$C(2026,1001) = 1026$\cdot$\ldots{}$\cdot$2026 tem o fator 2026 = 2$\cdot$1013, 1013 é primo e não divide 1001!, segue que 1013 | C(2026,1001). & tr(AB) = C(m+n,n-1) + C(m+n,n) + C(m+n+1,n+1) = C(m+n+2,n+1) = C(2026,1001) = 2026!/(1025!$\cdot$1001!), divisível por 1013. & igual & ambos \\
\bottomrule\end{longtable}%
\begin{itemize}%
\item%
\textbf{Demonstração (traço de A$\cdot$B divisível por 1013 para m = 2\^{}10, n = 10\^{}3))} Os dois chegam à mesma forma fechada, C(m+n+2,n+1) = C(2026,1001), e à mesma conclusão. Conferi tudo em Python: a identidade do traço vale para 1 $\leq$ m,n $\leq$ 24; as diagonais c11 = C(m+n,n-1), c22 = C(m+n,n) e c33 = C(m+n+1,n+1) do gabarito também estão certas; para m = 1024 e n = 1000 o traço calculado pela soma direta é igual a C(2026,1001), tem 609 algarismos, é $\equiv$ 0 (mod 1013) e não é divisível por 1013$^2$. O exemplo de conferência da resolução (m=2, n=1, T=10=C(5,2)) também está certo. A diferença está na justificativa final. O gabarito diz que 'o primo 1013 não divide 1025!$\cdot$1001!', o que é falso: 1013 $\leq$ 1025, então 1013 | 1025!. O argumento correto conta as potências de 1013: 2026! tem 1013$^2$ (dos fatores 1013 e 2026), 1025! tem 1013$^{1}$ e 1001! não tem, logo v\_1013 = 2-1-0 = 1. A resolução evita esse erro ao cancelar 1025! antes, usando 1001!$\cdot$C(2026,1001) = 1026$\cdot$\ldots{}$\cdot$2026, que contém apenas o múltiplo 2026 de 1013. Ela ainda confirma o resultado pelo teorema de Lucas e observa, corretamente, que a divisibilidade é exata de ordem 1. A demonstração da resolução está completa e correta. A conclusão do gabarito está certa, mas sua justificativa tem um passo falso.%
\end{itemize}%
\paragraph{Pontos de dúvida}%
\begin{itemize}%
\item%
\textit{(key)} A justificativa final do gabarito Poliedro tem um passo falso: diz que 'o primo 1013 não divide 1025!$\cdot$1001!', mas 1013 $\leq$ 1025, então 1013 divide 1025!. Na verdade 2026! tem 1013$^2$, 1025! tem 1013$^{1}$ e 1001! não tem fator 1013, de modo que v\_1013(C(2026,1001)) = 1. A conclusão continua certa, mas uma prova escrita como no gabarito seria logicamente inválida nesse ponto. Vale avisar o professor para não usar esse argumento como padrão de correção.%
\item%
\textit{(statement)} A hipótese m > n do enunciado não é necessária para a identidade tr(A$\cdot$B) = C(m+n+2,n+1), que vale para quaisquer inteiros positivos m e n, como a resolução observa. Não afeta a resposta.%
\end{itemize}%
\section*{Questão 8}%
\href{https://poliedroresolve.sistemapoliedro.com.br/vestibulares/ime/2027/ime-2-fase-dia-1-26-10-2026-2-fase-dia-1/dissertativa-8-matematica-2-fase-dia-1-ime-2027}{Resolução do gabarito}\par\smallskip%
\begin{longtable}{@{}p{0.06\textwidth} p{0.33\textwidth} p{0.33\textwidth} p{0.1\textwidth} p{0.08\textwidth}@{}}\toprule Item & Resolução & Gabarito & Veredito & Certo \\ \midrule
Única (demonstração da identidade) & Identidade demonstrada: S\_n = sen(n$\alpha$)/(sen(2$\alpha$)$\cdot$cos((n+1)$\alpha$)). Usa prostaférese, cos $\alpha$ + cos(2k+1)$\alpha$ = 2cos(k$\alpha$)cos((k+1)$\alpha$); telescopagem de sen $\alpha$/(cos(k$\alpha$)cos((k+1)$\alpha$)) = tg((k+1)$\alpha$) - tg(k$\alpha$); diferença de tangentes; e sen 2$\alpha$ = 2 sen $\alpha$ cos $\alpha$. Também trata o caso n = 0. & Demonstração: S\_n = sen(n$\alpha$)/(sen(2$\alpha$)$\cdot$cos((n+1)$\alpha$)), por prostaférese, telescopagem com tangentes e tg x - tg y = sen(x-y)/(cos x cos y). & igual & ambos \\
\bottomrule\end{longtable}%
\begin{itemize}%
\item%
\textbf{Única (demonstração da identidade))} O enunciado confere com o PDF original (p. 2, 8ª questão). A resolução segue o mesmo caminho do gabarito, com cada passo justificado. (1) Com P = $\alpha$ e Q = (2k+1)$\alpha$, o denominador vira 2cos((k+1)$\alpha$)cos(k$\alpha$), o que está correto. (2) Como sen $\alpha$ = sen((k+1)$\alpha$ - k$\alpha$), a identidade telescópica vale, e as divisões são permitidas porque k e k+1 estão em \{1,\ldots{},n+1\}. (3) A soma dá S\_n = (tg((n+1)$\alpha$) - tg $\alpha$)/(2 sen $\alpha$). (4) Essa diferença de tangentes é igual a sen(n$\alpha$)/(cos((n+1)$\alpha$)cos $\alpha$), o que leva ao resultado. Conferi numericamente em Python com (n, $\alpha$) = (1, $\pi$/6), (2, $\pi$/8), (3, $\pi$/3), (3, 0,71), (5, 0,37) e (7, 1,3): os dois lados coincidem. Todos os valores intermediários citados na resolução também estão corretos: 4,388842, -6,273816 e 0,987037 para $\alpha$ = 0,71; 0,76537 + 1,84776 = 2,61313 para $\alpha$ = $\pi$/8; -2 + 1 + 1 = 0 para $\alpha$ = $\pi$/3. A discussão sobre por que as hipóteses são necessárias também está certa: com $\alpha$ = $\pi$ e n = 1, S\_1 = -1/2, mas o lado direito fica 0/0; com $\alpha$ = $\pi$/4 e n = 1, os dois lados ficam indefinidos.%
\end{itemize}%
\paragraph{Pontos de dúvida}%
\begin{itemize}%
\item%
\textit{(statement)} O enunciado não diz se 0 $\in$ $\mathbb{N}$. A resolução trata n = 0 à parte (soma vazia, os dois lados valem 0), e o gabarito supõe implicitamente n $\geq$ 1. Isso não afeta a nota.%
\end{itemize}%
\section*{Questão 9}%
\href{https://poliedroresolve.sistemapoliedro.com.br/vestibulares/ime/2027/ime-2-fase-dia-1-26-10-2026-2-fase-dia-1/dissertativa-9-matematica-2-fase-dia-1-ime-2027}{Resolução do gabarito}\par\smallskip%
\begin{longtable}{@{}p{0.06\textwidth} p{0.33\textwidth} p{0.33\textwidth} p{0.1\textwidth} p{0.08\textwidth}@{}}\toprule Item & Resolução & Gabarito & Veredito & Certo \\ \midrule
Pontos de interseção do lugar geométrico de M com x = $\pm$c & (2$\surd$3, 1/7), (2$\surd$3, -1/7), (-2$\surd$3, 1/7), (-2$\surd$3, -1/7) $\approx$ ($\pm$3,46; $\pm$0,14) & (2$\surd$3, 1/7), (2$\surd$3, -1/7), (-2$\surd$3, 1/7), (-2$\surd$3, -1/7) & igual & ambos \\
Lugar geométrico de M (resultado intermediário) & 4x$^2$/49 + y$^2$ = 1 (elipse de semieixos 7/2 e 1), sem os pontos ($\pm$7/2, 0) & 4x$^2$/49 + y$^2$ = 1 & igual & ambos \\
\bottomrule\end{longtable}%
\begin{itemize}%
\item%
\textbf{Pontos de interseção do lugar geométrico de M com x = $\pm$c)} Refeito: a=4, b=2, c=$\surd$12=2$\surd$3. A normal em P=(x0,y0) tem direção (x0,4y0) (gradiente) e corta y=0 em Q=(3x0/4,0). Então M=(7x0/8, y0/2). Substituindo, chega{-}se a 4X$^2$/49+Y$^2$=1. Com X$^2$=12: Y$^2$=1-48/49=1/49, logo Y=$\pm$1/7. Conferi numericamente em Python com vários P parametrizados (o resultado deu exatamente 1 na equação do lugar). Os quatro pontos coincidem com os do gabarito.%
\item%
\textbf{Lugar geométrico de M (resultado intermediário))} A equação é a mesma. A resolução também exclui os pontos ($\pm$7/2, 0), que corresponderiam a P=($\pm$4,0); nesses pontos a normal é o próprio eixo x e Q fica indefinido. O gabarito não trata desse caso degenerado. A exclusão não mexe nas interseções pedidas, já que lá |y|=1/7$\neq$0.%
\end{itemize}%
\paragraph{Pontos de dúvida}%
\begin{itemize}%
\item%
\textit{(key)} O gabarito não trata os casos degenerados P=($\pm$4,0): nesses pontos a normal coincide com o eixo x e Q fica indefinido. A rigor, o lugar geométrico é a elipse 4x$^2$/49+y$^2$=1 sem os pontos ($\pm$7/2,0). A resolução registra isso e o gabarito não. Isso não afeta a resposta pedida.%
\item%
\textit{(statement)} O enunciado pede só as interseções com x=$\pm$c. O lugar geométrico entra apenas como passo intermediário. Também vale notar que c$^2$=12 é pouco menor que (7/2)$^2$=12,25: as retas cortam o lugar bem perto dos vértices, o que deixa |y|=1/7 pequeno. Esse valor está certo, não é erro de conta.%
\end{itemize}%
\section*{Questão 10}%
\href{https://poliedroresolve.sistemapoliedro.com.br/vestibulares/ime/2027/ime-2-fase-dia-1-26-10-2026-2-fase-dia-1/dissertativa-10-matematica-2-fase-dia-1-ime-2027}{Resolução do gabarito}\par\smallskip%
\begin{longtable}{@{}p{0.06\textwidth} p{0.33\textwidth} p{0.33\textwidth} p{0.1\textwidth} p{0.08\textwidth}@{}}\toprule Item & Resolução & Gabarito & Veredito & Certo \\ \midrule
a & V\_max = 192$\surd$3 cm$^3$ $\approx$ 332,55 cm$^3$ & V\_max = 192$\surd$3 cm$^3$ ($\approx$ 332,55 cm$^3$) & igual & ambos \\
b & Cubo: altura h = 4$\surd$3 cm $\approx$ 6,93 cm; base quadrada de lado 4$\surd$3 cm (diagonal da base 4$\surd$6 cm) & Cubo: altura e lados da base iguais a 4$\surd$3 cm & igual & ambos \\
\bottomrule\end{longtable}%
\begin{itemize}%
\item%
\textbf{a)} Os dois dão o mesmo valor e a mesma unidade. Refiz a conta: V $\leq$ h(72 - h$^2$/2), com máximo em h$^2$ = 48, o que dá V = 4$\surd$3$\cdot$48 = 192$\surd$3 $\approx$ 332,55 cm$^3$. A fatoração (h-4$\surd$3)$^2$(h+8$\surd$3) = h$^3$ - 144h + 384$\surd$3 foi conferida numericamente. Uma busca aleatória com prismas retangulares centrados encontrou $\approx$ 330,3, abaixo de 192$\surd$3, como esperado. A resolução vale para qualquer base quadrangular, convexa ou não (área $\leq$ $\frac12$$\cdot$d1$\cdot$d2 $\leq$ 2r$^2$), e para qualquer posição das bases. O gabarito supõe base retangular e os 8 vértices na superfície, mas chega ao mesmo valor.%
\item%
\textbf{b)} Os dois dão o mesmo resultado. Refiz a conta: r$^2$ = 36 - (h/2)$^2$ = 24, então o lado da base é r$\surd$2 = $\surd$48 = 4$\surd$3 cm, igual à altura (cubo de diagonal 4$\surd$3$\cdot$$\surd$3 = 12 = 2R). No gabarito aparece 'a = 4$\surd$(3 cm)', erro de digitação sem efeito no valor.%
\end{itemize}%
\paragraph{Pontos de dúvida}%
\begin{itemize}%
\item%
\textit{(key)} O gabarito parte de base retangular a$\times$b e de h = $\surd$(144 - a$^2$ - b$^2$), ou seja, supõe os 8 vértices na superfície da esfera. O enunciado só garante 4 vértices de uma face na superfície e diz que o prisma está contido na esfera, sem nada sobre a forma da base. A resposta não muda, porque o máximo sem restrição dentro da bola é o cubo inscrito, que cumpre a condição. Mas uma correção rigorosa poderia cobrar essa justificativa. A resolução conferida trata o caso geral (qualquer quadrilátero e qualquer posição das bases).%
\item%
\textit{(key)} Erro de digitação no gabarito: 'a = 4$\surd$(3 cm)' em vez de 'a = 4$\surd$3 cm'. Não muda o valor.%
\item%
\textit{(statement)} 'Contido em uma esfera' quer dizer contido na bola fechada. A frase sobre os 'quatro vértices de uma das faces' acaba não restringindo nada, porque o ótimo tem os 8 vértices na superfície. A leitura é natural e as duas soluções a adotam.%
\end{itemize}%
\end{document}